\documentclass[11pt]{article}
\usepackage[margin=1.08in]{geometry}
\usepackage{amsmath,amssymb,amsthm,mathtools}
\usepackage{booktabs}
\usepackage{enumitem}
\usepackage[hidelinks]{hyperref}
\usepackage{microtype}

\newtheorem{theorem}{Theorem}
\newtheorem{lemma}[theorem]{Lemma}
\newtheorem{proposition}[theorem]{Proposition}
\theoremstyle{remark}
\newtheorem{remark}[theorem]{Remark}

\newcommand{\tr}{\operatorname{tr}}
\newcommand{\eps}{\varepsilon}
\newcommand{\liminfT}{\liminf_{T\to\infty}}

\title{An Exact Pressure-Multiplicity Refinement for Simple Zeros of the Riemann Zeta Function}
\author{AMTOPA and contributors}
\date{August 12, 2026}

\begin{document}
\maketitle

\begin{abstract}
We record a short combinatorial refinement of the shifted-block averaging step in the positioned-pressure certificate of \texttt{sxuff/zeta-positioned-pressure}. Subject to the same imported analytic interface, the same certified seven-point inequality, and the same finite-dimensional Gram-profile estimate, the refinement replaces the coarse global pressure factor $m-1$ by the exact factor $m-6$. With the certified parameters of the predecessor, the resulting numerical lower bound is
\[
\liminfT \frac{N_0^s(T,2T)}{N(T,2T)}
\ge 0.673262375584978050386\ldots,
\]
so in particular it exceeds $0.6732623755$. This is a same-day research draft, not a peer-reviewed theorem and not a proof of the Riemann hypothesis.
\end{abstract}

\section{Statement and trust boundary}

Let $N(T,2T)$ denote the number of nontrivial zeros of $\zeta$ with ordinates in $(T,2T]$, counted with multiplicity, and let $N_0^s(T,2T)$ denote the number of simple zeros on the critical line in the same range. Write
\[
S=N_0^s(T,2T),\qquad N=N(T,2T).
\]

We import the analytic interface
\begin{equation}\label{eq:interface}
S\ge H_{\rm cert}N+\tr\Psi(M)-o(N),
\end{equation}
with
\[
H_{\rm cert}=\frac{672457041414544284}{10^{18}},
\]
and the certified seven-point inequality of the predecessor with
\[
\eps=\frac{51063}{10^7},\qquad b_r\ge 0,\qquad B:=\sum_{r=1}^6 b_r=\frac3{1150}.
\]
We also import the finite-dimensional block estimate
\begin{equation}\label{eq:hm-input}
\tr\Psi(G)\ge h_m(E),
\end{equation}
where
\[
h_m(E)=
\begin{cases}
E, & 0\le E\le \dfrac{m}{m-1},\\[1.2ex]
\dfrac{E}{m}+2\sqrt{\dfrac{m-1}{m}E}-1,
& E\ge \dfrac{m}{m-1}.
\end{cases}
\]
These inputs are not reproved here.

\begin{theorem}\label{thm:main}
Subject to the imported inputs above and the same normalization used in the predecessor's shifted-block argument,
\[
\liminfT \frac{N_0^s(T,2T)}{N(T,2T)}
\ge 0.673262375584978050386\ldots
>\frac{6732623755}{10^{10}}.
\]
\end{theorem}

\section{Exact local pressure multiplicity}

Consider a block of $m$ consecutive simple zeros, with normalized gaps
\[
g_1,\ldots,g_{m-1}.
\]
For each seven-point window starting at $t\in\{0,\ldots,m-7\}$, the local pressure contribution is
\[
\sum_{r=1}^{6} b_r g_{t+r}.
\]
After summing all $m-6$ local inequalities, define
\begin{equation}\label{eq:cj}
c_j:=\sum_{\substack{0\le t\le m-7\\1\le j-t\le 6}} b_{j-t},
\qquad
P_B:=\sum_{j=1}^{m-1}c_jg_j.
\end{equation}
The pair terms are handled exactly as in the predecessor, so with $E_B$ denoting the corresponding block off-diagonal energy,
\begin{equation}\label{eq:local-energy}
E_B+P_B\ge A_m,
\qquad
A_m:=\eps(m-6).
\end{equation}

\begin{lemma}[Exact pressure multiplicity]\label{lem:pressure}
The coefficients in \eqref{eq:cj} satisfy
\[
\sum_{j=1}^{m-1}c_j=(m-6)B.
\]
\end{lemma}

\begin{proof}
This is a direct double count of pairs consisting of a seven-point window and one of its six internal gap positions:
\begin{align*}
\sum_{j=1}^{m-1}c_j
&=\sum_{j=1}^{m-1}
  \sum_{\substack{0\le t\le m-7\\1\le j-t\le6}}b_{j-t}\\
&=\sum_{t=0}^{m-7}\sum_{r=1}^{6}b_r\\
&=(m-6)B.
\end{align*}
\end{proof}

\section{Retaining the exact pressure through the defect estimate}

Set
\[
R_m:=h_m(A_m),\qquad \eta_m:=\frac{R_m}{A_m}.
\]
Because $h_m$ is increasing and concave with $h_m(0)=0$, one has
\begin{equation}\label{eq:chord}
h_m(E)\ge \eta_m E\qquad(0\le E\le A_m).
\end{equation}

If $E_B\ge A_m$, monotonicity and \eqref{eq:hm-input} give
\[
\tr\Psi(G_B)\ge R_m-o(1).
\]
If $E_B<A_m$, then \eqref{eq:local-energy} gives $P_B\ge A_m-E_B$, and therefore by \eqref{eq:chord},
\begin{align*}
\tr\Psi(G_B)+\eta_mP_B
&\ge \eta_mE_B+\eta_m(A_m-E_B)-o(1)\\
&=R_m-o(1).
\end{align*}
Thus in all cases
\begin{equation}\label{eq:block}
\tr\Psi(G_B)+\eta_mP_B\ge R_m-o(1).
\end{equation}
The key point is that $P_B$ is retained in its exact position-dependent form; it is not replaced at this stage by the coarser $B\,\operatorname{span}(B)$.

\section{Shifted-block averaging}

Partition the ordered simple zeros into disjoint blocks of $m$ consecutive points, and consider all $m$ residue-class shifts of this partition. Up to $O(m)$ endpoint effects, a fixed global gap has the following behavior across the $m$ shifts:
\begin{enumerate}[label=(\roman*)]
\item it is a block boundary exactly once;
\item it occupies each internal position $j=1,\ldots,m-1$ exactly once.
\end{enumerate}
Consequently, by Lemma~\ref{lem:pressure}, its total pressure coefficient over all shifted partitions is exactly
\begin{equation}\label{eq:shift-multiplicity}
\sum_{j=1}^{m-1}c_j=(m-6)B.
\end{equation}
Averaging \eqref{eq:block} over the shifted partitions, using spectral pinching for the block defects and the same normalized total-gap relation as in the predecessor, gives
\begin{equation}\label{eq:global-defect}
\tr\Psi(M)
\ge \frac{R_m}{m}S
-\eta_mB\frac{m-6}{m}N
-o(N).
\end{equation}
For comparison, replacing each $c_j$ by the uniform upper bound $B$ before the shifted average yields the larger factor $B(m-1)/m$.

\section{Assembly and numerical specialization}

Substituting \eqref{eq:global-defect} into \eqref{eq:interface} gives
\[
S\ge H_{\rm cert}N+\frac{R_m}{m}S
-\eta_mB\frac{m-6}{m}N-o(N),
\]
and hence
\begin{equation}\label{eq:master}
\frac SN
\ge
\frac{mH_{\rm cert}-\eta_mB(m-6)}{m-R_m}-o(1).
\end{equation}
Since $A_m=\eps(m-6)$,
\[
\eta_mB(m-6)=\frac{R_mB}{\eps},
\]
so \eqref{eq:master} may also be written
\[
\frac SN
\ge
\frac{mH_{\rm cert}-R_mB/\eps}{m-R_m}-o(1).
\]

A high-precision scan of integer block lengths selects
\[
m=215.
\]
Then
\[
A_{215}=\frac{51063}{10^7}\cdot209=1.0672167,
\]
and
\[
R_{215}=1.06627687661618780546392584286702014367\ldots.
\]
Substitution into \eqref{eq:master} yields
\[
\frac SN
\ge
0.67326237558497805038619164619309578448\ldots-o(1).
\]
This proves Theorem~\ref{thm:main}, subject to the imported inputs.

\section{Reproducibility}

The repository contains two independent executable checks:
\begin{itemize}
\item \texttt{src/check\_multiplicity.py}, which verifies the exact rational identity in Lemma~\ref{lem:pressure} from the published pressure coefficients;
\item \texttt{src/check\_final\_bound.py}, which scans the block length and then recomputes the final bound with 100-digit interval arithmetic.
\end{itemize}
The GitHub Actions workflow runs both checks on every push and pull request.

\section{Research status and attribution}

This note contributes only the exact pressure-multiplicity retention through shifted-block averaging and the resulting arithmetic refinement. The analytic interface, the local seven-point interval certificate, the re-optimized window, and the sharp finite-$m$ block profile are inherited from the Anthropic / \texttt{trmdy} / \texttt{sxuff} lineage. If any imported input fails, the numerical bound in this note fails with it.

This manuscript should receive independent mathematical review before being cited as an established theorem.

\begin{thebibliography}{9}
\bibitem{sxuff}
\texttt{sxuff/zeta-positioned-pressure},
\url{https://github.com/sxuff/zeta-positioned-pressure}.

\bibitem{trmdy}
\texttt{trmdy/zeta-simple-zeros-673137},
\url{https://github.com/trmdy/zeta-simple-zeros-673137}.

\bibitem{anthropic}
Anthropic, research artifacts on simple zeros of the Riemann zeta function, 2026.
\end{thebibliography}

\end{document}
